\section{確率と統計的推論}
\label{D1_lesson16}

\subsection{授業内容}
\subsubsection{科目の中でのこのコマの位置づけ}
\begin{tabular}[t]{cccccc}
    記述統計[4] & 確率[3] & モデル[8] & 推測・推定[2] & 意思決定[0] & 実践[2]
\end{tabular}
\subsubsection{概要}
母集団に確率分布を仮定し，そこから得られる標本分布の特徴を利用して母数を推定する方法について学ぶ。標本の平均値が母平均に一致することを確認し，また標本の分散は母分散に一致しないことを確認する。ここでの推定方法はモーメントを用いたものであるが，確率分布を用いる方法もあり，その場合の``確率モデル''の考え方が必要なことを学ぶ。
\subsubsection{コマ主題細目(箇条書き)}
\begin{description}
\item[母集団と標本]推測統計学の基本的な概念である母集団と標本の関係を理解し，統計学的な推定(点推定，区間推定)の位置付けについて学ぶ。\footnote{本書では検定をモデル比較の一種として扱うので，ここでは言及しない。}用語としての母数(母平均，母分散，母比率など)と標本統計量(標本平均，標本分散，標本比率など)，および標本統計量の実現値といった用語の意味するところも慎重に理解しなければならない。
    \begin{flushright}
        $\to$ \citet{Yamada2004}のP.68--73
    \end{flushright}

\item[母集団分布とモデル]ここでは母集団に正規分布を仮定する``モデル''について考える。統計的な推測は基本的に不良設定問題であり，何らかの前提をおかないと解くことができない問題がほとんどである。心理学の場合は母集団分布を正規分布とする仮定が選ばれることが多く，この場合は標本統計量の分布も計算しやすくなる。
\begin{flushright}
    $\to$ \citet{Yamada2004}のP.74--79.
\end{flushright}

\item[標本分布と標準誤差]母集団が正規分布に従うとすると，標本統計量も正規分布に従うことが導出できる。ここでは標本統計量としての平均値を例にあげ，「標本統計量の分布」のことを標本分布と呼ぶこと，また標本分布の標準偏差を特に標準誤差と呼ぶことを理解する。たとえばサンプルサイズが大きくなったら，小さくなったらどのような値が得られるのかについて，具体的な例をみながら考えると理解が進む。
\begin{flushright}
    $\to$ 標本分布については\citet{Yamada2004}のP.90--97
\end{flushright}

\item[不偏推定量]母集団に正規分布を仮定し，標本平均の平均(期待値)が母平均に一致することを確認する。このことから，標本統計量が推定値として利用できることを確認する。ただし，分散の場合は一致しない(不偏性がない)ことをみる。分散については不偏分散を考えることでこの不一致を解決することができる。
\begin{flushright}
    $\to$ 不偏推定量については\citet{Yamada2004}のP.98--103に詳しい。
\end{flushright}

\end{description}
\subsubsection{キーワード}
\begin{itemize}
\item 母集団と標本
\item 標本分布
\item 点推定
\item 不偏推定量
\item 標本平均，標本分散
\end{itemize}
\subsection{授業情報}
\paragraph{コマの展開方法}
講義
\paragraph{標準シラバスにおける位置づけ}
科目番号5；心理学統計法；(2)統計に関する基礎的な知識；A 統計分析の基礎：母集団と標本
\subsubsection{予習・復習課題}
\paragraph{予習}
ここでは\citet{Yamada2004}が丁寧に解説されているので，事前に一読しておくことを強く勧める。
\paragraph{復習}
いよいよ記述統計をこえ，直接知り得ないものについて推論する段階に進む。基本的に解けない問題を解こうとするようなものであり，仮定やモデル，みなしている箇所について自覚的でなければならない。今回の講義の中で，何が仮定され，何が理論的に導出されたものなのかを区別できるかどうか，ノートを整理しておくと良いだろう。また推測統計学の基本的概念である，母集団と標本の関係と推定の手続きについては，\citet{Yamada2004}のテキストが最も丁寧で適切な記述をしているので，是非該当箇所を通読し，概念や用語を間違えないように注意されたし。
